\documentclass[11pt]{article}
\usepackage[margin=1in]{geometry}
\usepackage{booktabs}
\title{Convolutional Neural Network for Handwritten Digits}
\author{Devis W. Saputra}
\begin{document}
\maketitle

\begin{abstract}
This experiment compares a compact PyTorch CNN with flattened logistic regression on the scikit-learn handwritten-digits benchmark.
\end{abstract}

\section{Method}
Images are scaled to [0,1] and split 75/25 with stratification and seed 42. The CNN contains two convolution-pooling stages and a linear head, totaling 6,090 trainable parameters. Training uses Adam, cross-entropy loss, and a seeded DataLoader for deterministic batch ordering.

\section{Results}
\begin{tabular}{lrr}
\toprule
Model & Accuracy & Macro-F1 \\
\midrule
Logistic regression & 0.9622 & 0.9620 \\
CNN & 0.9733 & 0.9729 \\
\bottomrule
\end{tabular}

\section{Interpretation}
Unlike the MLP comparison, the CNN improves on its linear baseline. Preserving local two-dimensional image structure is useful in this setup.

\section{Limitations}
The benchmark is small. Stronger work would add repeated seeds, validation-based tuning, robustness checks, augmentation, and larger image datasets.

\section{Calculation audit and reproduction scope}
This PyTorch experiment compares a compact convolutional network with a flattened logistic-regression baseline on handwritten digits. The committed run reports 0.9733 accuracy for the CNN and 0.9622 for logistic regression after 14 epochs. The small gain illustrates a possible benefit from preserving spatial structure, while the single-split design and lack of an uncertainty estimate limit the strength of the architectural conclusion.

The recorded difference is about 1.11 percentage points on one split. This is not an uncertainty interval or evidence of broad image-recognition superiority. The MLP repository uses a different preprocessing setup and is not a direct cross-repository benchmark.

See the repository calculation guide for exact source paths. The full experiment was not independently rerun during this documentation review.
\end{document}
